\clearpage\begingroup\NoSourceHead\SourceLeaf{204}
\begin{center}
{\Large VERÖFFENTLICHUNGEN.}\\[1.8em]
{[I.]}\\[.8em]
\textbf{EINE LEICHTE METHODE, DEN OSTERSONNTAG ZU FINDEN.}\\[.7em]
\rule{.86\textwidth}{.35pt}\\[.2em]
{\small Astronomisches Jahrbuch für das Jahr 1814, Berlin 1811, S. 273.}\\[.2em]
\rule{.86\textwidth}{.35pt}
\end{center}
\noindent Im Julianischen Calender sey allemal $m=15$ und $n=6$.
\begin{center}
\begin{tabular}{@{}l r c r l l@{}}
Im Gregorian[ischen] von&1700&bis&1799&$m=23$&$n=3$.\\
&1800&&1899&$m=23$&$n=4$.
\end{tabular}
\end{center}
\begin{center}\small\setlength{\tabcolsep}{1.8pt}
\begin{tabular}{@{}cccccccccccl@{}}
Man&theile&das&vorgegebene&Jahr&durch&19&und&nenne&den&Rest&$a$,\\
»&»&»&»&»&»&4&»&»&»&»&$b$,\\
»&»&»&»&»&»&7&»&»&»&»&$c$,
\end{tabular}
\end{center}
\noindent Ferner [theile man]
\begin{center}\small\setlength{\tabcolsep}{3pt}
\begin{tabular}{@{}cclcccccc l@{}}
die&Anzahl&$(19a+m)$&durch&30&»&»&»&»&$d$,\\
»&»&$(2b+4c+6d+n)$&»&7&»&»&»&»&$e$.
\end{tabular}
\end{center}
\noindent So ist der \textls[120]{Ostersonntag} in beiden Calendern der $22+d+e$ März*).
\begingroup\renewcommand{\thefootnote}{}\footnotetext{*) Wenn im gregor[ianischen] Calender die Rechnung Ostern am 26st. April giebt, setzt man allemal den 19t. und wenn sie den 25st. bringt, den 18t.}\endgroup
\begin{center}Z. B. für 1812.\end{center}
\[
\begin{array}{r@{\;}l}
&1812\\[-8pt]19)&\rule{2em}{.35pt}\\[-3pt]&\text{rest. }7=a
\end{array}
\qquad
\begin{array}{r@{\;}l}
&1812\\[-8pt]4)&\rule{2em}{.35pt}\\[-3pt]&\text{rest. }0=b
\end{array}
\qquad
\begin{array}{r@{\;}l}
&1812\\[-8pt]7)&\rule{2em}{.35pt}\\[-3pt]&\text{rest. }6=c.
\end{array}
\]
\begin{center}Für den Jul[ianischen] Calender.\end{center}
\[
\begin{array}{r@{\;}l}
(19\mathbin{.}7)+15=&148\\[-6pt]30)&\rule{2em}{.35pt}\\[-3pt]&\text{rest. }28=d
\end{array}
\qquad
\begin{array}{r@{\;}l}
(2\mathbin{.}0)+(4\mathbin{.}6)+(6\mathbin{.}28)+6=&198\\[-6pt]7)&\rule{2em}{.35pt}\\[-3pt]&\text{rest. }2=e
\end{array}
\]
\[
22+28+2\text{ März}=52\text{ März}=21\text{ April der Ostersonntag.}
\]
\clearpage\SourceLeaf{205}\HistoricalHead{200}{L}{CHRONOLOGIE.\quad VERÖFFENTLICHUNGEN.}
\begin{center}Für den Gregor[ianischen] Calender.\end{center}
\[
\begin{array}{r@{\;}l}
(19\mathbin{.}7)+23=&156\\[-6pt]30)&\rule{2em}{.35pt}\\[-3pt]&\text{rest. }6=d
\end{array}
\qquad
\begin{array}{r@{\;}l}
(2\mathbin{.}0)+(4\mathbin{.}6)+(6\mathbin{.}6)+4=&64\\[-6pt]7)&\rule{2em}{.35pt}\\[-3pt]&\text{rest. }1=e
\end{array}
\]
\[
22+6+1\text{ März}=29\text{ März der Ostersonntag.}
\]
\begin{center}\textbf{BEMERKUNG.}\end{center}
\begingroup\small
\noindent\hspace*{\parindent}Bei den Angaben dieser Notiz hat \textsc{Gauss} nicht berücksichtigt, daß in den Jahren 1700 bis 1899 Ostern im gregorianischen Kalender zweimal auf den 25. April fällt, nämlich 1734 und 1886. Für diese zwei Jahre ergibt die Rechnung den 25. April, und Ostern ist auch tatsächlich an diesem Tage zu begehen, da die Zahl $11m+11$ durch 30 dividiert, einen Rest liefert, der größer als 19 ist; nur sofern $11m+11$, durch 30 dividiert, einen Rest gibt, der kleiner als 19 ist, wird Ostern, wenn aus der Rechnung der 25. April folgt, nicht an diesem Tage, sondern am 18. April begangen, wie \textsc{Gauss} in seiner ersten Veröffentlichung*) richtig angibt.\begingroup\renewcommand{\thefootnote}{}\footnotetext{*) \textsc{Gauss}, \textit{Berechnung des Osterfestes}, Monatliche Correspondenz zur Beförderung der Erd- und Himmels-Kunde, herausgeg. vom Freiherrn v. \textsc{Zach} 2, August 1800, S. 121, Werke VI, S. 73 ff., vgl. S. 79 unter II; die für die Jahre 1700 bis 1899 gültige Zahl $m=23$ gehört auch nicht zu den a. a. O. S. 79, Z. 8 v. u. verzeichneten acht Ausnahmewerten. Auch \textsc{Gauss}' Angabe in seiner zweiten Veröffentlichung, \textit{Noch etwas über die Bestimmung des Osterfestes}, Braunschweigisches Magazin 1807 September 12, Werke VI, S. 85, Z. 15 v. u., nach der für die Verlegung des Osterfestes vom 25. auf den 18. April erforderlich ist, daß der Rest der Division der Jahreszahl durch 19 mindestens gleich 11 wird, ist völlig korrekt; diese Bedingung trifft offenbar für die Jahre 1734 und 1886 nicht zu.}\endgroup\ In ihr führt er sogar 1734 und 1886 als Jahre mit dem absolut spätesten Osterdatum an (Werke VI, S. 74, vorletzte und letzte Zeile). Daß Ostern unter Umständen am 25. April zu feiern ist, hat übrigens auch der Freiherr v. \textsc{Zach}**) bei seiner ersten Wiedergabe der \textsc{Gauss}schen Osterformel nicht beachtet.\begingroup\renewcommand{\thefootnote}{}\footnotetext{**) \textit{Tables abrégées et portatives de la lune} par le Baron \textsc{de Zach}, Florence 1809. Die Osterformel findet man a. a. O. S. 67, auf S. XIII wird \textsc{Gauss} als ihr Urheber genannt.}\endgroup\ Trotzdem er sich an \textsc{Gauss}' erste Publikation in der Beziehung anschließt, ist nach seinen Angaben Ostern \textls[120]{stets} am 18. April zu begehen, wenn sich durch die Rechnung der 25. April ergibt. Auch \textsc{Delambre} hat in seinem \textit{Abrégé d'Astronomie}, Paris 1813, S. 647 sowie in seiner \textit{Astronomie théorique et pratique} III, Paris 1814, S. 712 die unrichtige Angabe, daß Ostern im gregorianischen Kalender nie auf den 25. April fallen kann. Seine irrigen Angaben hat \textsc{Delambre} später***) verbessert.\begingroup\renewcommand{\thefootnote}{}\footnotetext{***) \textsc{Delambre}, \textit{Formules pour calculer la Lettre Dominicale, le Nombre d'Or, l'Epacte et la fête de Pâques, pour une année Grégorienne ou Julienne quelconque}, Connaissance des tems pour l'an 1817, Paris 1815, S. 307 ff., insbesondere S. 315; vergl. auch \textsc{Delambre}, \textit{Histoire de l'Astronomie moderne} I, Paris 1821, S. 24.}\endgroup\ Auch v. \textsc{Zach} hat$\dagger$) bei einer erneuten Wiedergabe der \textsc{Gauss}schen Osterformel richtige Angaben über den 25. April als Osterdatum im gregorianischen Kalender gemacht.\begingroup\renewcommand{\thefootnote}{}\footnotetext{$\dagger$) Baron \textsc{de Zach}, \textit{Notes}, Correspondance astronomique I, Genua 1818, S. 568.}\endgroup
\begin{flushright}\textsc{Loewy.}\end{flushright}
\endgroup
\begin{center}\rule{.24\textwidth}{.35pt}\end{center}
\clearpage\endgroup
